% Generated from locale/tr/content/first-order-logic/sequent-calculus/proving-things.tex; do not hand-edit.


\iftag{FOL}
      {\olfileid[tr]{fol}{seq}{pro}}
      {\olfileid[tr]{pl}{seq}{pro}}

\olsection{\usetoken{P}{derivation} Örnekleri}

\begin{ex}
$!A \land !B \Sequent !A$ sekantı için bir $\Log{LK}$-!!{derivation} verin.

İstenen son sekantı !!{derivation}{}in en altına yazarak başlarız.
\begin{prooftree}
\AxiomC{}
\UnaryInf$!A\land !B \fCenter !A$
\end{prooftree}
Sonra, alt sekantı bu biçimde olan çıkarımın hangi türden olabileceğini
belirlememiz gerekir. Bu bir yapısal kural olabilir; ancak işe mantıksal bir
kural arayarak başlamak yerinde olur. Alt sekantta görünen tek mantıksal
bağlaç $\land$ olduğundan bir $\land$ kuralı ararız. $\land$ simgesi öncülde
bulunduğuna göre \LeftR{\land} kuralına bakarız.
\begin{prooftree}
\AxiomC{}
\RightLabel{\LeftR{\land}}
\UnaryInf$!A\land !B \fCenter !A$
\end{prooftree}
\LeftR{\land} çıkarımının üst sekantı için iki seçenek vardır: $!A
\Sequent !A$ ya da $!B \Sequent !A$. $!A \Sequent !A$ açıkça bir başlangıç
sekantıdır (bu iyi bir şeydir); buna karşılık $!B \Sequent !A$ genel olarak
türetilemez. Üst sekantı dolduralım:
\begin{prooftree}
\Axiom$!A \fCenter !A$
\RightLabel{\LeftR{\land}}
\UnaryInf$!A\land !B \fCenter !A$
\end{prooftree}
Artık $!A \land !B \Sequent !A$ sekantının doğru bir
$\Log{LK}$-!!{derivation}{}ini elde ettik.
\end{ex}

\begin{ex}
$\lnot !A \lor !B \Sequent !A \lif !B$ sekantı için bir
$\Log{LK}$-!!{derivation} verin.

İstenen son sekantı !!{derivation}{}in en altına yazarak başlayın.
\begin{prooftree}
\AxiomC{}
\UnaryInf$\lnot !A \lor !B \fCenter !A \lif !B$
\end{prooftree}
Bu son sekantı verebilecek bir mantıksal kural bulmak için son sekanttaki
mantıksal bağlaçlara bakarız: $\lnot$, $\lor$ ve $\lif$. Şimdilik yalnızca
$\lor$ ile $\lif$ bizi ilgilendirir; çünkü bunlar son sekanttaki
!!{main operator}s olup !!{sentence}s içinde yer alır; $\lnot$ ise başka bir bağlacın
kapsamındadır ve onu sonra ele alacağız. Buna göre son çıkarım için mantıksal
kural seçeneklerimiz \LeftR{\lor} ile \RightR{\lif} kurallarıdır.
Aslında ikisinden birini seçebiliriz; fakat iki dala ayrılmayı ertelememizi
sağladığı için \RightR{\lif} kuralını seçelim. Bu alt sekantı verebilecek
\RightR{\lif} çıkarımlarının biçimine göre şunu elde etmeliyiz:
\begin{prooftree}
\AxiomC{}
\UnaryInf$ !A, \lnot !A \lor !B \fCenter !B $
\RightLabel{\RightR{\lif}} \UnaryInf$ \lnot !A \lor !B \fCenter !A \lif !B $
\end{prooftree}

$\lnot !A \lor !B$ ifadesini öncülün dışına taşırsak \LeftR{\lor} kuralını
uygulayabiliriz. Şemaya göre bunun aşağıdaki iki üst sekanta ayrılması gerekir:
\begin{prooftree}
\AxiomC{}
\UnaryInf$\lnot !A, !A \fCenter !B$
\AxiomC{}
\UnaryInf$!B, !A \fCenter !B$
\RightLabel{\LeftR{\lor}}
\BinaryInf$ \lnot !A \lor !B, !A \fCenter !B $
\RightLabel{\LeftR{\Exchange}}
\UnaryInf$ !A, \lnot !A \lor !B \fCenter !B $
\RightLabel{\RightR{\lif}}
\UnaryInf$ \lnot !A \lor !B \fCenter !A \lif !B $
\end{prooftree}
Başlangıç sekantlarına doğru yukarı ilerlemeye çalıştığımızı hatırlayın;
neredeyse ulaştık!{} Sağ dal, bir başlangıç sekantından yalnızca bir
zayıflatma ve bir yer değiştirme uzaktadır; bunları uyguladığımızda bu dal
tamamlanır:
\begin{prooftree}
\AxiomC{}
\UnaryInf$\lnot !A, !A \fCenter !B$
\Axiom$!B \fCenter !B$
\RightLabel{\LeftR{\Weakening}}
\UnaryInf$!A, !B \fCenter !B$
\RightLabel{\LeftR{\Exchange}}
\UnaryInf$!B, !A \fCenter !B$
\RightLabel{\LeftR{\lor}}
\BinaryInf$\lnot !A \lor !B, !A \fCenter !B $
\RightLabel{\LeftR{\Exchange}}
\UnaryInf$ !A, \lnot !A \lor !B \fCenter !B $
\RightLabel{\RightR{\lif}}
\UnaryInf$ \lnot !A \lor !B \fCenter !A \lif !B $
\end{prooftree}

Şimdi sol dala baktığımızda herhangi bir !!{sentence} içinde kalan tek mantıksal
bağlacın, öncülü oluşturan !!{sentence}s arasındaki $\lnot$ simgesi olduğunu görürüz;
dolayısıyla \LeftR{\lnot} kuralının bir örneğini arıyoruz.
\begin{prooftree}
\AxiomC{}
\UnaryInf$ !A \fCenter !B, !A$
\RightLabel{\LeftR{\lnot}}
\UnaryInf$\lnot !A, !A \fCenter !B$
\Axiom$!B \fCenter !B$
\RightLabel{\LeftR{\Weakening}}
\UnaryInf$!A, !B \fCenter !B$
\RightLabel{\LeftR{\Exchange}}
\UnaryInf$!B, !A \fCenter !B$
\RightLabel{\LeftR{\lor}}
\BinaryInf$\lnot !A \lor !B, !A \fCenter !B $
\RightLabel{\LeftR{\Exchange}}
\UnaryInf$ !A, \lnot !A \lor !B \fCenter !B $
\RightLabel{\RightR{\lif}}
\UnaryInf$ \lnot !A \lor !B \fCenter !A \lif !B $
\end{prooftree}
Sağ dalı tamamladığımız gibi bu sol dalı da tamamlamak için yalnızca bir
zayıflatma ve bir yer değiştirme gerekir.
\begin{prooftree}
\Axiom$!A \fCenter !A$
\RightLabel{\RightR{\Weakening}}
\UnaryInf$ !A \fCenter !A, !B$
\RightLabel{\RightR{\Exchange}}
\UnaryInf$ !A \fCenter !B, !A$
\RightLabel{\LeftR{\lnot}}
\UnaryInf$\lnot !A, !A \fCenter !B$
\Axiom$!B \fCenter !B$
\RightLabel{\LeftR{\Weakening}}
\UnaryInf$!A, !B \fCenter !B$
\RightLabel{\LeftR{\Exchange}}
\UnaryInf$!B, !A \fCenter !B$
\RightLabel{\LeftR{\lor}}
\BinaryInf$\lnot !A \lor !B, !A \fCenter !B $
\RightLabel{\LeftR{\Exchange}}
\UnaryInf$ !A, \lnot !A \lor !B \fCenter !B $
\RightLabel{\RightR{\lif}}
\UnaryInf$ \lnot !A \lor !B \fCenter !A \lif !B $
\end{prooftree}
\end{ex}

\begin{ex}
$\lnot !A \lor \lnot !B \Sequent \lnot (!A \land !B)$ sekantının bir
$\Log{LK}$-!!{derivation}{}ini verin.

Yukarıdaki teknikleri kullanarak istenen son sekantı en alta yazmakla
başlarız.
\begin{prooftree}
\AxiomC{}
\UnaryInf$ \lnot !A \lor \lnot !B \fCenter \lnot (!A \land !B) $
\end{prooftree}
Son sekanttaki !!{sentence}s bakımından ana bağlaçlar $\lor$ ile $\lnot$ simgeleridir.
Burada \LeftR{\lor} ya da \RightR{\lnot} kuralını uygulamak işe yarayabilir;
ancak bir süre daha iki dala ayrılmaktan kaçındığı için \RightR{\lnot}
kuralıyla başlarız:
\begin{prooftree}
\AxiomC{}
\UnaryInf$!A \land !B, \lnot !A \lor \lnot !B \fCenter $
\RightLabel{\RightR{\lnot}}
\UnaryInf$\lnot !A \lor \lnot !B \fCenter \lnot (!A \land !B)$
\end{prooftree}
Şimdi \LeftR{\land} ya da \LeftR{\lor} kuralına bakabiliriz.
\LeftR{\land} kuralını uygularsak ne olacağını görelim: $!A,
\lnot !A \lor \lnot !B \Sequent \quad$ veya $!B,\lnot !A \lor \lnot !B
\Sequent \quad$ üst sekantlarından birini seçebiliriz. Bu !!{derivation}, $!A$
ile $!B$ bakımından simetrik olduğundan ilkini seçelim:
\begin{prooftree}
\AxiomC{}
\UnaryInf$!A, \lnot !A \lor \lnot !B \fCenter $
\RightLabel{\LeftR{\land}}
\UnaryInf$!A \land !B, \lnot !A \lor \lnot !B \fCenter $
\RightLabel{\RightR{\lnot}}
\UnaryInf$\lnot !A \lor \lnot !B \fCenter \lnot (!A \land !B)$
\end{prooftree}
Bu !!{derivation}{}i doldurmaya devam edince bir sorunla karşılaşırız:
\begin{prooftree}
\Axiom$!A \fCenter !A$
\RightLabel{\LeftR{\lnot}}
\UnaryInf$ \lnot !A, !A \fCenter$
\AxiomC{}
\RightLabel{?}
\UnaryInf$!A \fCenter !B$
\RightLabel{\LeftR{\lnot}}
\UnaryInf$ \lnot !B, !A \fCenter$
\RightLabel{\LeftR{\lor}}
\BinaryInf$\lnot !A \lor \lnot !B, !A \fCenter $
\RightLabel{\LeftR{\Exchange}}
\UnaryInf$!A, \lnot !A \lor \lnot !B \fCenter $
\RightLabel{\LeftR{\land}}
\UnaryInf$!A \land !B, \lnot !A \lor \lnot !B \fCenter $
\RightLabel{\RightR{\lnot}}
\UnaryInf$\lnot !A \lor \lnot !B \fCenter \lnot (!A \land !B)$
\end{prooftree}
Sağ dalın tepesi daha fazla indirgenemez ve yapısal çıkarımlarla bir başlangıç
sekantına dönüştürülemez; dolayısıyla izlenecek doğru yol bu değildir. Öyleyse
yukarıda \LeftR{\land} kuralını uygulamak hataydı. Bir önceki aşamaya dönüp
onun yerine \LeftR{\lor} kuralını uygularsak şunu elde ederiz:
\begin{prooftree}
\AxiomC{}
\UnaryInf$\lnot !A, !A \land !B \fCenter $

\AxiomC{}
\UnaryInf$\lnot !B, !A \land !B \fCenter $

\RightLabel{\LeftR{\lor}}
\BinaryInf$\lnot !A \lor \lnot !B, !A \land !B \fCenter $
\RightLabel{\LeftR{\Exchange}}
\UnaryInf$!A \land !B, \lnot !A \lor \lnot !B \fCenter $
\RightLabel{\RightR{\lnot}}
\UnaryInf$\lnot !A \lor \lnot !B \fCenter \lnot (!A \land !B)$
\end{prooftree}
Her dalı daha önce yaptığımız biçimde tamamlayınca şu sonucu elde ederiz:
\begin{prooftree}
\Axiom$ !A \fCenter!A$
\RightLabel{\LeftR{\land}}
\UnaryInf$!A \land !B \fCenter !A$
\RightLabel{\LeftR{\lnot}}
\UnaryInf$\lnot !A, !A \land !B \fCenter $

\Axiom$ !B \fCenter !B$
\RightLabel{\LeftR{\land}}
\UnaryInf$!A \land !B \fCenter !B$
\RightLabel{\LeftR{\lnot}}
\UnaryInf$\lnot !B, !A \land !B \fCenter $

\RightLabel{\LeftR{\lor}}
\BinaryInf$\lnot !A \lor \lnot !B, !A \land !B \fCenter $
\RightLabel{\LeftR{\Exchange}}
\UnaryInf$!A \land !B, \lnot !A \lor \lnot !B \fCenter $
\RightLabel{\RightR{\lnot}}
\UnaryInf$\lnot !A \lor \lnot !B \fCenter \lnot (!A \land !B)$
\end{prooftree}
(Bu adımlarda $\land$ kurallarını $\lnot$ kurallarından daha aşağıda uygulayıp
yine doğru bir !!{derivation} elde edebilirdik.)
\end{ex}

\begin{ex}
Şimdiye kadar daraltma kuralını kullanmadık; ancak bu kural bazen gereklidir.
Bunun gerçekleştiği bir örneğe bakalım. $\quad \Sequent !A \lor \lnot !A$
sekantını ispatlamak istediğimizi varsayalım. $\RightR{\lor}$ kuralını geriye
doğru uygulamak bize şu !!{derivation}s arasından birini verir:
\begin{prooftree}
\AxiomC{}
\UnaryInf$ \fCenter !A$
\RightLabel{\RightR{\lor}}
\UnaryInf$ \fCenter !A \lor \lnot !A$
\DisplayProof\qquad\bottomAlignProof
\AxiomC{}
\UnaryInf$!A \fCenter $
\RightLabel{\RightR{\lnot}}
\UnaryInf$ \fCenter \lnot !A$
\RightLabel{\RightR{\lor}}
\UnaryInf$ \fCenter !A \lor \lnot !A$
\end{prooftree}
Elbette bunların hiçbiri bir başlangıç sekantında son bulmaz. İşin püf noktası,
daraltma kuralının, bir !!a{sentence} söz konusu olduğunda onun iki kopyasını tek kopyada birleştirmemizi
sağladığını fark etmektir. Bir ispat ararken, yani aşağıdan yukarıya giderken,
öncülde $!A \lor \lnot !A$ ifadesinin bir kopyasını tutabiliriz; örneğin:
\begin{prooftree}
\AxiomC{}
\UnaryInf$ \fCenter !A \lor \lnot !A, !A$
\RightLabel{\RightR{\lor}}
\UnaryInf$ \fCenter !A \lor \lnot !A, !A \lor \lnot !A$
\RightLabel{\RightR{\Contraction}}
\UnaryInf$ \fCenter !A \lor \lnot !A$
\end{prooftree}
Artık $\RightR{\lor}$ kuralını ikinci kez uygulayıp $\lnot !A$'yı da elde
edebiliriz; böylece eksiksiz bir !!{derivation}{}e ulaşırız.
\begin{prooftree}
\Axiom$!A \fCenter !A$
\RightLabel{\RightR{\lnot}}
\UnaryInf$\fCenter !A, \lnot !A$
\RightLabel{\RightR{\lor}}
\UnaryInf$\fCenter !A, !A \lor \lnot !A$
\RightLabel{\RightR{\Exchange}}
\UnaryInf$ \fCenter !A \lor \lnot !A, !A$
\RightLabel{\RightR{\lor}}
\UnaryInf$ \fCenter !A \lor \lnot !A, !A \lor \lnot !A$
\RightLabel{\RightR{\Contraction}}
\UnaryInf$ \fCenter !A \lor \lnot !A$
\end{prooftree}
\end{ex}

\begin{prob}
Aşağıdaki sekantlar için !!{derivation}s verin:
\begin{enumerate}
\item $!A \land (!B \land !C) \Sequent (!A \land !B) \land !C$.
\item $!A \lor (!B \lor !C) \Sequent (!A \lor !B) \lor !C$.
\item $!A \lif (!B \lif !C) \Sequent !B \lif (!A \lif !C)$.
\item $!A \Sequent \lnot\lnot !A$.
\end{enumerate}
\end{prob}

\begin{prob}
Aşağıdaki sekantlar için !!{derivation}s verin:
\begin{enumerate}
\item $(!A \lor !B) \lif !C \Sequent !A \lif !C$.
\item $(!A \lif !C) \land (!B \lif !C) \Sequent (!A \lor !B) \lif !C$.
\item $\Sequent \lnot(!A \land \lnot !A)$.
\item $!B \lif !A \Sequent \lnot !A \lif \lnot !B$.
\item $\Sequent (!A \lif \lnot !A) \lif \lnot !A$.
\item $\Sequent \lnot(!A \lif !B) \lif \lnot !B$.
\item $!A \lif !C \Sequent \lnot (!A \land \lnot !C)$.
\item $!A \land \lnot !C \Sequent \lnot (!A \lif !C)$.
\item $!A \lor !B, \lnot !B \Sequent !A$.
\item $\lnot !A \lor \lnot !B \Sequent \lnot(!A \land !B)$.
\item $\Sequent (\lnot !A \land \lnot !B) \lif\lnot(!A \lor !B)$.
\item $\Sequent \lnot(!A \lor !B) \lif (\lnot !A \land \lnot !B)$.
\end{enumerate}
\end{prob}

\begin{prob}
Aşağıdaki sekantlar için !!{derivation}s verin:
\begin{enumerate}
\item $\lnot(!A \lif !B) \Sequent !A$.
\item $\lnot(!A \land !B) \Sequent \lnot !A \lor \lnot !B$.
\item $!A \lif !B \Sequent \lnot !A \lor !B$.
\item $\Sequent \lnot \lnot !A \lif !A$.
\item $!A \lif !B, \lnot !A \lif !B \Sequent !B$.
\item $(!A \land !B) \lif !C \Sequent (!A \lif !C) \lor (!B \lif !C)$.
\item $(!A \lif !B) \lif !A \Sequent !A$.
\item $\Sequent (!A \lif !B) \lor (!B \lif !C)$.
\end{enumerate}
(Bunların tümü $\RightR{\Contraction}$ kuralını gerektirir.)
\end{prob}
